\section{Recovery Uncertainty With Unknown Onset}
\label{app:unknown-onset-bridge}
\newcommand{\BridgeMomentOne}{$+3.31$}
\newcommand{\BridgeMomentOneCI}{[$2.01$, $4.82$]}
\newcommand{\BridgeMixtureOne}{$-0.29$}
\newcommand{\BridgeMixtureOneCI}{[$-1.22$, $0.57$]}
\newcommand{\BridgeMomentTwo}{$+3.38$}
\newcommand{\BridgeMomentTwoCI}{[$1.91$, $4.93$]}
\newcommand{\BridgeMixtureTwo}{$+0.96$}
\newcommand{\BridgeMixtureTwoCI}{[$0.07$, $1.91$]}

We compared three representations of a two-component recovery distribution
on a linear swing model: a Gaussian with the shared within-component
covariance, a Gaussian matching both moments, and the full mixture.
A mean-only control omits state uncertainty altogether.
The likelihood integrates a bounded Gaussian severity and marginalizes
over a finite onset grid. A separate set of normal windows calibrates each
arm to a 5\% false alarm probability per window. The observation model
uses two angles and an approximate branch flow.

In the base setting, adding between-component dispersion increased
joint detection and source coverage by \BridgeMomentOne{} percentage
points (exploratory paired 95\% interval \BridgeMomentOneCI{}).
The change was \BridgeMomentTwo{} points (\BridgeMomentTwoCI{}) on
a second seed. Retaining the full mixture rather than matching its
moments changed coverage by \BridgeMixtureOne{} and \BridgeMixtureTwo{}
points (intervals \BridgeMixtureOneCI{} and \BridgeMixtureTwoCI{}). Normal false alarm rates in test differed by arm.
These simulations use a specified initial recovery law, evaluate complete
windows rather than continuous alarms and have no topology changes.
They motivate a control that preserves the second moment in a future
bank. They do not isolate the effect of mixture retention in the
IEEE 39-bus system.\EOD
